\section{Project Risk Management}

Risk management is the process through which potential problems are identified, assessed, and prepared for before they seriously affect project success \cite{risk1}. In the proposed AI-Driven Lab Exam Monitoring and Management System, risk management is especially important because the project involves technical integration, real-time workflow assumptions, multiple stakeholder expectations, and strict academic deadlines. Risks do not need to become actual failures if they are recognized early and managed in a disciplined way.

\subsection{Risk Identification}

Several categories of risk may affect the project. Technical risks include integration difficulties between modules, instability in monitoring features, or data-handling complexity \cite{facerec1}. Schedule risks include delays in design completion, incomplete sprint outcomes, or documentation falling behind implementation. Academic risks include misunderstanding stakeholder needs, insufficient justification in the report, or weak alignment between written chapters and the evolving prototype. Resource risks may include limited access to realistic testing conditions, hardware constraints, or tool-related limitations.

\subsection{Risk Analysis}

Risk analysis helps the team determine which risks deserve the most attention \cite{risk1}. A risk is not important only because it is possible; it becomes important when its likelihood combines with meaningful impact. For example, a small wording correction in a chapter may be likely but low impact, whereas instability in a monitoring feature during late-stage integration may be both moderately likely and highly disruptive. The team therefore evaluates risks in terms of both probability and seriousness.


\begin{table}[H]
\centering
\caption{Project Risk Rating Matrix}
\label{tab:risk_matrix}
\renewcommand{\arraystretch}{1.2}
\begin{tabular}{>{\raggedright\arraybackslash}p{7cm}|ccc}
\hline\hline
\textbf{Risk} & \textbf{Likelihood} & \textbf{Impact} & \textbf{Priority} \\ \hline
Face recognition accuracy under varied lab conditions & Medium & High & High \\
Incomplete module integration before deadline         & Medium & High & High \\
Real-time monitoring latency or system overload       & Low    & High & Medium \\
Workstation or hardware unavailability in lab         & Low    & High & Medium \\
Third-party library or API deprecation                & Low    & Medium & Low \\
Data privacy or student record exposure               & Very Low & High & Medium \\
Team member unavailability during critical sprint     & Low    & Medium & Low \\ \hline\hline
\end{tabular}
\end{table}

\subsection{Risk Mitigation Plan}

Risk mitigation involves planning responses before problems grow. Requirement instability can be reduced through early scope clarification and regular supervisor review. Integration risk can be reduced through modular development and early testing of interfaces between key components. Documentation delay can be reduced by writing chapters incrementally throughout the project rather than postponing them. Testing constraints can be reduced through scenario-based validation even when full real-world deployment is not possible. Time pressure can be reduced through milestone tracking and careful attention to critical-path activities.

\subsection{Risk Monitoring and Control}

Risk management is not a one-time exercise performed only at the proposal stage. Risks must be monitored throughout the semester as the project evolves. New risks may appear when modules become more detailed or when feedback changes priorities. Therefore, the team should review risk status regularly, especially after each sprint or milestone. Continuous risk monitoring supports more stable project execution and reduces the likelihood that late-stage surprises will undermine the final deliverable.

